\section{CONFIG-009: GITCLAW\_ENV Only Partially Supported}
\label{sec:config-009}

\subsection*{Metadata}
\begin{tabular}{ll}
\textbf{Issue ID} & CONFIG-009 \\
\textbf{Severity} & Medium \\
\textbf{Category} & Configuration \\
\textbf{Branch} & \branchref{fix/CONFIG-009-gitclaw-env-support} \\
\textbf{Changes} & \branchdiff{fix/CONFIG-009-gitclaw-env-support} \\
\end{tabular}

\subsection*{Executive Summary}
The \texttt{GITCLAW\_ENV} environment variable is used in \texttt{src/config.ts:45} to determine which config preset to load, but it is only referenced in \texttt{loadEnvConfig()} and nowhere else. Other parts of the system that could benefit from environment-aware behavior---different log levels, error verbosity, model fallbacks---do not check \texttt{GITCLAW\_ENV}. Additionally, there is no validation that the environment name is one of the expected values (\texttt{development}, \texttt{staging}, \texttt{production}); any string is accepted and silently ignored.

The fix adds environment validation with a warning for unknown values and introduces environment-aware variables in the entry point for downstream use.

\subsection*{Root Cause Analysis}

The original code only used \texttt{GITCLAW\_ENV} in one place:

\begin{lstlisting}[language=TypeScript]
// src/config.ts:43-55 — Before fix
export async function loadEnvConfig(
    agentDir: string, env?: string
): Promise<EnvConfig> {
    const configDir = join(agentDir, "config");
    const envName = env || process.env.GITCLAW_ENV;  // Only usage in codebase

    const base = await loadYamlFile(join(configDir, "default.yaml"));
    if (envName) {
        const envOverride = await loadYamlFile(join(configDir, envName + ".yaml"));
        return deepMerge(base, envOverride) as EnvConfig;
    }
    return base as EnvConfig;
}
\end{lstlisting}

There was no validation of the environment name---any string passed as \texttt{GITCLAW\_ENV} would be used as a filename without asserting it was a known environment.

\subsection*{Fix Implementation}

The fix adds a set of valid environments and validation with a warning:

\begin{lstlisting}[language=TypeScript]
// src/config.ts — After fix
const VALID_ENVS = new Set(["development", "staging", "production", "test"]);

export async function loadEnvConfig(
    agentDir: string, env?: string
): Promise<EnvConfig> {
    const configDir = join(agentDir, "config");
    const envName = env || process.env.GITCLAW_ENV;

    if (envName && !VALID_ENVS.has(envName)) {
        console.warn(
            "[config] Unknown environment \"" + envName + "\". " +
            "Valid values: " + [...VALID_ENVS].join(", "));
    }

    const base = await loadYamlFile(join(configDir, "default.yaml"));
    if (envName) {
        const envOverride = await loadYamlFile(join(configDir, envName + ".yaml"));
        return deepMerge(base, envOverride) as EnvConfig;
    }
    return base as EnvConfig;
}
\end{lstlisting}

Additionally, environment-aware variables were introduced in the entry point:

\begin{lstlisting}[language=TypeScript]
// src/index.ts — After fix
const _envName = process.env.GITCLAW_ENV || "production";
const _isDev = _envName === "development";
\end{lstlisting}

These variables can be used throughout the system to adjust behavior based on the environment, such as enabling verbose logging in development or full stack traces.

\subsection*{Affected Files}
\begin{itemize}
    \item \srcfile{src/config.ts} --- added \texttt{VALID\_ENVS} set and validation
    \item \srcfile{src/index.ts} --- added \texttt{\_envName} and \texttt{\_isDev} constants
    \item \srcfile{src/\_\_tests\_\_/config-009.test.ts} --- new environment validation tests
\end{itemize}

\subsection*{Verification}
\begin{lstlisting}[language=TypeScript]
// config-009.test.ts
const VALID_ENVS = new Set(["development", "staging", "production", "test"]);

function validateEnv(name: string | undefined): boolean {
    if (!name) return true;
    if (!VALID_ENVS.has(name)) {
        console.warn("Unknown environment \"" + name + "\"");
        return false;
    }
    return true;
}

test("valid environment names are accepted", () => {
    strictEqual(validateEnv("development"), true);
    strictEqual(validateEnv("staging"), true);
    strictEqual(validateEnv("production"), true);
    strictEqual(validateEnv("test"), true);
});

test("invalid environment name is rejected", () => {
    strictEqual(validateEnv("invalid-env"), false);
});

test("undefined environment is accepted (no env set)", () => {
    strictEqual(validateEnv(undefined), true);
});
\end{lstlisting}
